\documentclass{article}
\usepackage{amsmath,amssymb,mathtools}
\usepackage[margin=1in]{geometry}
\allowdisplaybreaks
\usepackage[UTF8]{ctex}
\begin{document}
\title{最优水流的流向混合整数分段线性松弛算法}
\date{}
\maketitle

\section*{集合与数据}
该有向网络为$G=(\mathcal N,\mathcal L)$。其节点划分方式为$\mathcal N=\mathcal D\mathbin{\dot\cup}\mathcal R\mathbin{\dot\cup}\mathcal T$，包括需求节点（包括零需求节点）、水库以及储水池。其链接划分方式为$\mathcal L=\mathcal A\mathbin{\dot\cup}\mathcal V\mathbin{\dot\cup}\mathcal P$，包括管道、阀门以及固定速度的泵。
\[
\delta_i^+=\{(i,j)\in\mathcal L\},\qquad
\delta_i^-=\{(j,i)\in\mathcal L\},\qquad
\deg_i^+=|\delta_i^+|,\quad \deg_i^-=|\delta_i^-|.
\]
对于整数$K\ge1$，设$\mathcal K=\{1,\ldots,K\}$和$\widetilde{\mathcal K}=\{1,\ldots,K+1\}$。索引$k$处的时间值为$t(k)$，而$k$到$k+1$之间的间隔长度为$\Delta t^k$。

以下所有系数均为固定实例数据。供应头部边界$\underline h_i^k,\overline h_i^k$适用于$i\in\mathcal N$、$k\in\widetilde{\mathcal K}$；规定的水库头部边界为$H_i^k$；需求值$\overline q_i^k$适用于$i\in\mathcal D$、$k\in\mathcal K$；储水池节点流量边界$\underline q_i^k,\overline q_i^k$适用于$i\in\mathcal T$、$k\in\mathcal K$。正节点流量表示向网络中的注入量，负节点流量表示抽取量。储水池数据包括直径$D_i$、底部高程$B_i^k$以及初始容积$\widetilde v_i$，其中底部高程低于最低水位。

请提供$\underline q_{ij}^k,\overline q_{ij}^k$的签名链接流界限，以及非负的方向性下限$\underline q_{ij}^{k+},\underline q_{ij}^{k-}$。泵浦的开启状态下限$\underline q_{ij}^{k+}$为正数。方向性上限为……
\begin{equation}
\overline q_{ij}^{k+}=\max\{0,\overline q_{ij}^k\},\qquad
\overline q_{ij}^{k-}=\max\{0,-\underline q_{ij}^k\},\qquad
\underline q_{ij}^{ks}\ge0\quad(s\in\{+,-\}).
\tag{D1}
\end{equation}
泵的反向流限制为零。对于每条管道，供给长度$L_{ij}>0$以及每单位长度的阻力$r_{ij}$；管道的 exponent 为$\alpha=1.852$。泵的系数满足$a_{ij}>0$、$b_{ij}<0$、$c_{ij}>0$，并定义了源头的严格凹形泵曲线类别。供电成本系数$\lambda_{ij}^{\mathrm{cost},k},\mu_{ij}^k$，最小泵启停持续时间$\tau^{\mathrm{on}},\tau^{\mathrm{off}}$，以及激活限制$N_{ij}$。

对于每一个管道、运行时间以及方向$s\in\{+,-\}$，请提供有序的分割点数组。
\[
\mathcal Q_{ij}^{ks}=(\widehat q_{ijp}^{ks})_{p=1}^{n_{ij}^{ks}},\qquad n_{ij}^{ks}=|\mathcal Q_{ij}^{ks}|.
\]
对于每一个泵和运行时间，都需要提供$\mathcal Q_{ij}^{k+}=(\widehat q_{ijp}^{k+})_{p=1}^{n_{ij}^{k+}}$的值。所显示的端点及间隔索引值，需要用到$n_{ij}^{ks}\ge2$来表示每个数组的情况。所有断点功率和切线系数都必须是有限实数。它们的数值，比如网络规模和其他参数值，都是实例输入。

\section*{变量和域}
真正的变量是头部$h_i^k$（$i\in\mathcal N$, $k\in\widetilde{\mathcal K}$）、节点流$q_i^k$（$i\in\mathcal N$, $k\in\mathcal K$）、带符号的链接流$q_{ij}^k$（$ (i,j)\in\mathcal L$, $k\in\mathcal K$），以及储罐容积$v_i^k$（$i\in\mathcal T$, $k\in\widetilde{\mathcal K}$）。这些变量仅具有下面所述的界限和等式关系。链接流使用带符号的实数域。方向性流、管道头部差异以及泵的输出量则使用不同的域。
\begin{align}
&q_{ij}^{ks}\ge0 &&((i,j)\in\mathcal L,\ k\in\mathcal K,\ s\in\{+,-\}),\
&\Delta h_{ij}^{ks}\ge0 &&((i,j)\in\mathcal A,\ k\in\mathcal K,\ s\in\{+,-\}),\
&g_{ij}^k\ge0 &&((i,j)\in\mathcal P,\ k\in\mathcal K).
\end{align}
二进制域是
\begin{align}
&y_{ij}^k\in\{0,1\} &&((i,j)\in\mathcal L,\ k\in\mathcal K),\
&z_{ij}^k\in\{0,1\} &&((i,j)\in\mathcal V\cup\mathcal P,\ k\in\mathcal K),\
&z_{ij}^{\mathrm{on},k},z_{ij}^{\mathrm{off},k}\in\{0,1\} &&((i,j)\in\mathcal P,\ k\in\mathcal K).
\end{align}
对于管道符号，每个$k\in\mathcal K$和$s\in\{+,-\}$对应的分段变量都满足以下条件：
\[
0\le\lambda_{ijp}^{ks}\le1\quad(p=1,\ldots,n_{ij}^{ks}),\qquad
x_{ijp}^{ks}\in\{0,1\}\quad(p=2,\ldots,n_{ij}^{ks}).
\]
对于泵类设备，同样的规则适用于$\lambda_{ijp}^{k+}$和$x_{ijp}^{k+}$，其中$n_{ij}^{k+}$代替了$n_{ij}^{ks}$。这些$\lambda$变量与固定的电力成本系数$\lambda^{\mathrm{cost}}$是不同的。

使用仿射方向别名
\begin{equation}
y_{ij}^{k+}:=y_{ij}^k,\qquad y_{ij}^{k-}:=1-y_{ij}^k.
\tag{D2}
\end{equation}

\section*{目标}
\begin{equation}
\min\ \sum_{k\in\mathcal K}\sum_{(i,j)\in\mathcal P}
\left(\lambda_{ij}^{\mathrm{cost},k}q_{ij}^k+\mu_{ij}^kz_{ij}^k\right).
\tag{14}
\end{equation}

\section*{节点与时间限制}
\begin{align}
\underline h_i^k\le h_i^k\le\overline h_i^k
&\quad(i\in\mathcal N,\ k\in\widetilde{\mathcal K}), \tag{1}\\
h_i^k=H_i^k
&\quad(i\in\mathcal R,\ k\in\widetilde{\mathcal K}), \tag{R1}\\
q_i^k=\overline q_i^k
&\quad(i\in\mathcal D,\ k\in\mathcal K), \tag{N1}\\
q_i^k\ge0
&\quad(i\in\mathcal R,\ k\in\mathcal K), \tag{N2}\\
\underline q_i^k\le q_i^k\le\overline q_i^k
&\quad(i\in\mathcal T,\ k\in\mathcal K), \tag{N3}\\
v_i^k=\frac{\pi}{4}D_i^2(h_i^k-B_i^k)
&\quad(i\in\mathcal T,\ k\in\widetilde{\mathcal K}), \tag{2}\\
v_i^{k+1}=v_i^k-\Delta t^kq_i^k
&\quad(i\in\mathcal T,\ k\in\mathcal K), \tag{3}\\
v_i^1=\widetilde v_i,\qquad v_i^1\le v_i^{K+1}
&\quad(i\in\mathcal T). \tag{12}
\end{align}
\begin{equation}
\sum_{(i,j)\in\delta_i^+}q_{ij}^k-
\sum_{(j,i)\in\delta_i^-}q_{ji}^k=q_i^k
\quad(i\in\mathcal N,\ k\in\mathcal K).
\tag{11}
\end{equation}

\section*{阀门和泵}
对于每一个$(i,j)\in\mathcal V$、$k\in\mathcal K$，
\begin{align}
\underline q_{ij}^kz_{ij}^k\le q_{ij}^k\le\overline q_{ij}^kz_{ij}^k, \tag{6}\\
(1-z_{ij}^k)(\underline h_i^k-\overline h_j^k)
\le h_i^k-h_j^k\le
(1-z_{ij}^k)(\overline h_i^k-\underline h_j^k). \tag{7}
\end{align}
对于每一个$(i,j)\in\mathcal P$、$k\in\mathcal K$，
\begin{align}
\underline q_{ij}^{k+}z_{ij}^k\le q_{ij}^k\le\overline q_{ij}^kz_{ij}^k, \tag{8}\\
q_{ij}^{k-}=0, \tag{P1}\\
h_i^k-h_j^k+g_{ij}^k\le(1-z_{ij}^k)(\overline h_i^k-\underline h_j^k), \tag{10a}\\
h_i^k-h_j^k+g_{ij}^k\ge(1-z_{ij}^k)(\underline h_i^k-\overline h_j^k). \tag{10b}
\end{align}
处于关闭状态的增益为零，这一点可以通过非负增益限制以及下方的泵浦包络约束来体现。

对于每一个$(i,j)\in\mathcal P$，施加…
\begin{align}
z_{ij}^k-z_{ij}^{k-1}\le z_{ij}^{\mathrm{on},k}
&\quad(k\in\mathcal K\setminus\{1\}), \tag{13b}\\
z_{ij}^{\mathrm{on},k}\le z_{ij}^{k'}
&\quad(k,k'\in\mathcal K:\ t(k)\le t(k')\le t(k)+\tau^{\mathrm{on}}), \tag{13c}\\
z_{ij}^{k-1}-z_{ij}^k\le z_{ij}^{\mathrm{off},k}
&\quad(k\in\mathcal K\setminus\{1\}), \tag{13d}\\
z_{ij}^{k'}\le1-z_{ij}^{\mathrm{off},k}
&\quad(k,k'\in\mathcal K:\ t(k)\le t(k')\le t(k)+\tau^{\mathrm{off}}), \tag{13e}\\
\sum_{k\in\mathcal K}z_{ij}^{\mathrm{on},k}\le N_{ij}. & \tag{13f}
\end{align}

\section*{方向性表述}
对于每一个$(i,j)\in\mathcal L$、$k\in\mathcal K$，
\begin{align}
q_{ij}^k=q_{ij}^{k+}-q_{ij}^{k-}, \tag{15}\\
\underline q_{ij}^{k+}y_{ij}^k\le q_{ij}^{k+}\le\overline q_{ij}^{k+}y_{ij}^k, \tag{16b}\\
\underline q_{ij}^{k-}(1-y_{ij}^k)\le q_{ij}^{k-}\le\overline q_{ij}^{k-}(1-y_{ij}^k). \tag{16c}
\end{align}
对于每一个$(i,j)\in\mathcal A$、$k\in\mathcal K$，
\begin{align}
\Delta h_{ij}^{k+}-\Delta h_{ij}^{k-}=h_i^k-h_j^k, \tag{17a}\\
0\le\Delta h_{ij}^{k+}\le y_{ij}^k\max\{0,\overline h_i^k-\underline h_j^k\}, \tag{17b}\\
0\le\Delta h_{ij}^{k-}\le(1-y_{ij}^k)\max\{0,\overline h_j^k-\underline h_i^k\}. \tag{17c}
\end{align}

\section*{管状信封}
对于每一个$(i,j)\in\mathcal A$、$k\in\mathcal K$、$s\in\{+,-\}$，都需施加相应的切线不等式。
\begin{equation}
r_{ij}\left[\widehat q^{\alpha}y_{ij}^{ks}
+\alpha\widehat q^{\alpha-1}(q_{ij}^{ks}-\widehat qy_{ij}^{ks})\right]
\le\frac{\Delta h_{ij}^{ks}}{L_{ij}}
\quad(\widehat q\in\mathcal Q_{ij}^{ks}),
\tag{18}
\end{equation}
以及分段信封和邻接约束
\begin{align}
\sum_{p=1}^{n_{ij}^{ks}}r_{ij}(\widehat q_{ijp}^{ks})^\alpha\lambda_{ijp}^{ks}
&\ge\frac{\Delta h_{ij}^{ks}}{L_{ij}}, \tag{22}\\
q_{ij}^{ks}&=\sum_{p=1}^{n_{ij}^{ks}}\widehat q_{ijp}^{ks}\lambda_{ijp}^{ks}, \tag{23}\\
\sum_{p=1}^{n_{ij}^{ks}}\lambda_{ijp}^{ks}=y_{ij}^{ks},\qquad
\sum_{p=2}^{n_{ij}^{ks}}x_{ijp}^{ks}&=y_{ij}^{ks}, \tag{24a}\\
\lambda_{ijp}^{ks}&\le x_{ijp}^{ks}+x_{ij,p+1}^{ks}
\quad(p=2,\ldots,n_{ij}^{ks}-1), \tag{24b}\\
\lambda_{ij,1}^{ks}\le x_{ij,2}^{ks},\qquad
\lambda_{ij,n_{ij}^{ks}}^{ks}&\le x_{ij,n_{ij}^{ks}}^{ks}. \tag{24c}
\end{align}

\section*{将信封中的内容抽出。}
对于每一个$(i,j)\in\mathcal P$、$k\in\mathcal K$，都施加…
\begin{equation}
g_{ij}^k\le[a_{ij}+b_{ij}\widehat q^{c_{ij}}]z_{ij}^k
+b_{ij}c_{ij}\widehat q^{c_{ij}-1}(q_{ij}^{k+}-\widehat qz_{ij}^k)
\quad(\widehat q\in\mathcal Q_{ij}^{k+}),
\tag{20}
\end{equation}
\begin{align}
g_{ij}^k&\ge\sum_{p=1}^{n_{ij}^{k+}}
\left(a_{ij}+b_{ij}(\widehat q_{ijp}^{k+})^{c_{ij}}\right)\lambda_{ijp}^{k+}, \tag{25}\\
q_{ij}^{k+}&=\sum_{p=1}^{n_{ij}^{k+}}\widehat q_{ijp}^{k+}\lambda_{ijp}^{k+}, \tag{26}\\
\sum_{p=1}^{n_{ij}^{k+}}\lambda_{ijp}^{k+}=z_{ij}^k,\qquad
\sum_{p=2}^{n_{ij}^{k+}}x_{ijp}^{k+}&=z_{ij}^k, \tag{27a}\\
\lambda_{ijp}^{k+}&\le x_{ijp}^{k+}+x_{ij,p+1}^{k+}
\quad(p=2,\ldots,n_{ij}^{k+}-1), \tag{27b}\\
\lambda_{ij,1}^{k+}\le x_{ij,2}^{k+},\qquad
\lambda_{ij,n_{ij}^{k+}}^{k+}&\le x_{ij,n_{ij}^{k+}}^{k+}. \tag{27c}
\end{align}

\section*{基于方向的合法不等式}
对于每一个$k\in\mathcal K$，施加…
\begin{align}
\sum_{(i,j)\in\delta_i^+}y_{ij}^k&\ge1
&&\left(i\in\mathcal R\cup\{i'\in\mathcal D:\overline q_{i'}^k>0\}\right), \tag{32}\\
\sum_{(j,i)\in\delta_i^-}y_{ji}^k+
\sum_{(i,j)\in\delta_i^+}(1-y_{ij}^k)&\ge1
&&(i\in\mathcal D:\overline q_i^k<0), \tag{33}\\
\sum_{(j,i)\in\delta_i^-}y_{ji}^k-
\sum_{(i,j)\in\delta_i^+}y_{ij}^k&=0
&&(i\in\mathcal D:\overline q_i^k=0,\ \deg_i^+=\deg_i^-=1), \tag{34a}\\
\sum_{(j,i)\in\delta_i^-}y_{ji}^k+
\sum_{(i,j)\in\delta_i^+}y_{ij}^k&=1
&&\left(\begin{array}{l}i\in\mathcal D:\overline q_i^k=0,\\(\deg_i^+,\deg_i^-)\in\{(2,0),(0,2)\}\end{array}\right). \tag{34b}
\end{align}
\end{document}
